% ===== tr-07 块5：第 5 章 pp.16--19（原文页 16--19）=====

% [原文 p.16]
对任意 $p>1$，把方程乘以 $\phi|\phi|^{p-2}$ 并分部积分，
\[
\int_X|\nabla|\phi|^{p/2}|_\omega^2\omega^n \le \frac{C(n,\|f\|_{L^\infty})p^2}{p-1}\int_X|\phi|^{p-1}\omega^n.
\]
关键在于右端的较弱范数控制了左端的较强范数。这与 Poincaré 不等式
\[
\int_X u^2\omega^n\le C\int_X|du|_\omega^2\omega^n,\qquad \int_X u=0,
\]
以及 \textbf{Sobolev 不等式}
\[
\int_X|u|^{\frac{2n}{n-1}}\omega^n\le C\Big(\int_X|du|_\omega^2\omega^n+\int_X|u|^2\omega^n\Big)
\]
的方向恰好相反。二者合力给出 $\|\phi\|_{L^2_1}\le C$，并且可以归纳地估计大 $p$ 时的 $\|\phi\|_{L^p}$，而该估计对 $p$ 是一致成立的。取极限 $p\to\infty$，就得到界 $\|\phi\|_{L^\infty}\le C(X,\|f\|_{L^\infty})$。

% [Tutor Note: Moser 迭代的关键是“弱范数控制强范数”（右端的 $L^{p-1}$ 控制左端的梯度 $L^2$），这与 Poincaré 不等式、Sobolev 不等式“用强范数控制弱范数”的方向相反；Sobolev 指数正是 $\frac{2n}{n-1}$，不可误写为 $\frac{2n}{n+1}$。]

第二步是估计 $\|dd^c\phi\|_{C^0}$（称为 $C^{1,\bar1}$-估计），即证明 $\omega$ 与 $\omega_\phi$ 之间的一致等价。出发点是局部计算给出的微分不等式（采用分析学家的 Laplacian 约定）
\[
\Delta_{\omega_\phi}\log\operatorname{Tr}_\omega\omega_\phi\ge -C\operatorname{Tr}_{\omega_\phi}\omega-C.
\]
减去恒等式
\[
\Delta_{\omega_\phi}\phi = n-\operatorname{Tr}_{\omega_\phi}\omega
\]
的足够大的倍数，便得到微分不等式
\[
\Delta_{\omega_\phi}(\log\operatorname{Tr}_\omega\omega_\phi-C\phi)\ge \operatorname{Tr}_{\omega_\phi}\omega-C'.
\]
利用 $\phi$ 的先验界，再应用\textbf{极大值原理}即可推出 $\operatorname{Tr}_\omega\omega_\phi\le C$。

% [Tutor Note: 恒等式 $\Delta_{\omega_\phi}\phi=n-\operatorname{Tr}_{\omega_\phi}\omega$ 是这套论证的枢纽：它把 $\operatorname{Tr}_{\omega_\phi}\omega$ 换成 $n-\Delta_{\omega_\phi}\phi$，从而可以“减去 $C\phi$”把坏项吸收到左边、再对 $\log\operatorname{Tr}_\omega\omega_\phi-C\phi$ 用极大值原理。抄写时注意 $\operatorname{Tr}_{\omega_\phi}\omega$ 的下标顺序（相对 $\omega_\phi$ 取迹）。]

\begin{remark}{5.4}
这一 $C^{1,\bar1}$-估计论证具有\textbf{整体性}：如果我们只知道 (5.9) 在 $\C^n$ 中一个标准单位球上成立、且 $\phi$ 有给定的 $C^0$-界，那么 $C^{1,\bar1}$-界是可以出现反例的，解可能发展出奇点。
\end{remark}

第三步是得到高阶估计。标准椭圆理论表明，只要 $\phi$ 有 $C^{2,\alpha}$ 界就足够了。Evans--Krylov 理论弥合了 $C^{1,\bar1}$-界与 $C^{2,\alpha}$-界之间的空隙（细节参看 [66] 第 2 章第 4 节）。这一论证具有\textbf{局部性}，并且关键地用到：$\log\det(\partial_i\partial_{\bar j}\phi)$ 是矩阵 $(\partial_i\partial_{\bar j}\phi)$ 的\emph{凹}函数；主要工具是 Harnack 不等式。其几何洞见在于：高阶正则性是 Kähler 流形局部欧氏本性的体现。

% [Tutor Note: 第三步与前两步的性质相反：$C^{1,\bar1}$-估计必须整体（见注 5.4），而 Evans--Krylov 的高阶估计是局部的，其前提是已有 $C^{1,\bar1}$ 界。原文 $C^{1,\bar1}$ 中的上标是 $\bar1$（横杠在 1 上），不是普通的 $1,1$。]

Yau 的证明策略影响深远，渗透到绝大多数关于 CY 度量的工作之中。然而上面这份简短的梗概凸显了它难以移植到 SYZ 猜想情形的两个原因：

% [原文 p.17]
\begin{itemize}
\item 经历大复结构极限的 CY 度量高度退化，以致 Sobolev 常数变得过大，从而使 Moser 迭代技巧无法给出关于势的有效 $C^0$-估计。
\item $C^{1,\bar1}$ 界具有整体性，因此若要只在一般区域上获得关于 Calabi--Yau 度量的有用估计，极大值原理的设计必须在整体上成立。这在高度退化的流形上是很困难的，因为那里的黎曼曲率无法整体地一致有界。
\end{itemize}

由于这些原因，我们处理 SYZ 猜想的方法与 Yau 的证明有重大偏离。势的估计转而使用复势理论，而一般区域中的度量估计则使用椭圆偏微分方程理论中 Savin 的一个深刻定理，从而绕开了 Yau 的估计。

% [Tutor Note: 这两条原因正好对应 Yau 证明的两个关键步骤（Moser 迭代的 $C^0$-估计、整体极大值原理的 $C^{1,\bar1}$-估计）；作者指出在退化族上二者都失效，因此本文改为“复势理论 + Savin 定理”。这一对应关系是理解第 4 节结构的关键。]

\subsection{复势理论}

复势理论研究 Kähler 势的各种弱概念，以及它们对代数几何与偏微分方程的更广泛推论。可以把它们看作黎曼曲面上的势理论在多个复变量情形下的推广。一些常见的主题包括：
\begin{itemize}
\item 势空间的弱紧性理论。它适用于 Kähler 几何中的变分方法。
\item 势理论与代数几何之间的联系，经由 Hodge 理论、$\bar\partial$-算子、相交理论等等。这为双有理几何提供了超越方法。
\item 复 Monge--Ampère 方程的弱形式，以及在体积密度上的弱假设之下、或环境复簇存在奇点时的先验势估计。这对于研究典范度量的奇点形成是有用的。
\end{itemize}
与 Yau 对 Calabi 猜想的证明（其技巧是典型的椭圆偏微分方程技巧）不同，复势理论更接近于复分析，并且常常给出更强的结果。为了建立直觉，我们将先回顾文献中的标准版本，然后陈述我们的统一版本------后者是我们处理 SYZ 猜想的基础。

\subsubsection{Skoda 不等式}

$\C^n$ 中坐标球上的上半连续 $L^1$-函数 $\phi$ 称为多重次调和的（psh），若它限制到每条复直线上都满足次中值不等式；这蕴含 $dd^c\phi\ge0$。基本的直觉是：

% [原文 p.18]
一般维数下 psh 函数的正则性性质类似于黎曼曲面上的次调和函数。这一点由 Skoda 不等式的局部版本所刻画：

\begin{theorem}{5.4.2}
（参看 [79, Thm 3.1]）若 $\phi$ 在 $B_2\subset\C^n$ 上是 psh，且关于标准欧氏度量 $\omega_E$ 有 $\int_{B_2}|\phi|\omega_E^n\le1$，则存在仅依赖于维数的常数 $\alpha$、$C$，使得
\[
\log\int_{B_1}e^{-\alpha\phi}\omega_E^n\le C.
\]
\end{theorem}

\begin{remark}{5.5}
Skoda 不等式可以与 $\R^{2n}$ 单位球上的次调和函数（$n>1$）的情形相对照：在那里指数可积性远远超出了可期望的范围。
\end{remark}

% [Tutor Note: 注意球的半径下标：假设中积分在 $B_2$ 上、结论中积分在 $B_1$ 上，这是一个“缩小半径”的局部估计，不能把两处都写成 $B_1$ 或 $B_2$。]

在紧 Kähler 流形 $(Y,\omega)$ 上，若上半连续 $L^1$-函数 $\phi$ 与 $\omega$ 的局部势之和是 psh 的，即 $\omega_\phi=\omega+dd^c\phi\ge0$，则称 $\phi\in PSH(Y,\omega)$。这是 Kähler 势的推广概念。Skoda 不等式的标准整体版本是：

\begin{theorem}{5.4.3}
[73] 在固定的 $(X,\omega)$ 上，存在只依赖于 $X,\omega$ 的正数常数 $\alpha$、$C$，使得
\[
\int_X e^{-\alpha\phi}\omega^n\le C,\qquad \forall\phi\in PSH(X,\omega)\ \text{且}\ \sup\phi=0.
\]
\end{theorem}

在我们的应用中，我们需要处理大复结构极限附近 Calabi--Yau 流形的极化代数退化 $\pi:X\to S\setminus\{0\}$，如第 5.3 节所述。设 $\omega_{FS}$ 是 $(X,c_1(L))$ 上一个固定的 Fubini--Study 度量，由 $L$ 的高次幂的截面所给出的射影嵌入诱导；并用
\[
\omega_{FS,t}=\frac{1}{|\log|t||}\omega_{FS}|_{X_t}
\]
来定义 $X_t$ 上类 $\frac{1}{|\log|t||}c_1(L)$ 中一族背景度量。归一化因子 $\frac{1}{|\log|t||}$ 是为了保证：两个不同的 Fubini--Study 参考度量的选取所差的势，其 $C^0$ 范数为与足够小的 $t$ 无关的 $O(1)$ 阶。回忆 $d\mu_t$ 是归一化 Calabi--Yau 测度。

我们把 Skoda 不等式改造成如下统一版本 [51]：

\begin{theorem}{5.4.4}
（统一 Skoda 估计）存在与 $t$（$0<|t|\ll1$）无关的一致正数常数 $\alpha$、$A$，使得对于归一化 Calabi--Yau 测度 $d\mu_t$，
\[
\int_{X_t}e^{-\alpha u}d\mu_t\le A,\qquad \forall u\in PSH(X_t,\omega_{FS,t})\ \text{且}\ \sup_{X_t}u=0.
\]
\end{theorem}

% [Tutor Note: 定理 5.4.4 与定理 5.4.3 的差别在于：测度从 $\omega^n$ 换成了 $d\mu_t$，度量从固定 $\omega$ 换成了随 $t$ 变化的 $\omega_{FS,t}$，而常数 $\alpha,A$ 与 $t$ 无关------这正是“统一（uniform）”的含义。]

其证明用大量看起来像 $\C^n$ 中标准球的区域覆盖 $X_t$。一个初等但略显微妙的试验函数构造使得人们可以估计局部势的 $L^1$ 范数。然后对每个小区域应用 Skoda 不等式的局部版本，并对所有区域求和。

% [原文 p.19]
\subsection{势函数的估计}

复势理论中的一个基本问题是对 Kähler 势作估计。Kolodziej 开创了一种方法，用以获得如下效果。他的定理的基本版本在固定的环境 Kähler 流形 $(Y,\omega)$ 上工作，处理归一化为 $\sup_Y\phi=0$ 的 Kähler 势 $\phi\in PSH(Y,\omega)$。
\begin{itemize}
\item （“势估计”）若 $\omega_\phi^n$ 的体积密度具有某种可积性控制，例如一个 $L^p$-界
\begin{equation}
\int_Y\Big|\frac{\omega_\phi^n}{\omega^n}\Big|^p\le C,\qquad p>1, \tag{5.10}
\end{equation}
则 $\phi$ 有只依赖于 $(X,\omega),n,p,C$ 的 $L^\infty$-界 [45][20][24]。事实上这可以改进为 $\phi$ 的 $C^\alpha$ Hölder 界 [47]。（注意 $p=1$ 是不够的，因为 $L^1$-界 $\int_Y\omega_\phi^n\le\int_Y\omega^n$ 是自动成立的。）
\item （“$L^1$-稳定性估计”）设 $\phi,\psi\in PSH(Y,\omega)$ 都满足 $L^p$ 体积密度可积性控制 (5.10)，且归一化 $\sup_Y\phi=\sup_Y\psi=0$。若
\begin{enumerate}
\item 或者 $\phi,\psi$ 在 $L^1$-意义下彼此接近，即 $\int_Y|\psi-\phi|\omega^n\ll1$（“$L^1$-势稳定性”）；
\item 或者 $\omega_\phi^n$ 与 $\omega_\psi^n$ 的体积密度在 $L^1$-意义下彼此接近，即全变差 $\int_Y|\omega_\phi^n-\omega_\psi^n|\ll1$（“$L^1$-体积稳定性”）；
\end{enumerate}
那么 $|\phi-\psi|$ 在 $L^\infty$-意义下很小，且有定量估计 [46]。
\end{itemize}

% [Tutor Note: (5.10) 中的被积函数是体积密度之比 $\omega_\phi^n/\omega^n$ 的 $p$ 次幂；$p>1$ 这一条件是本质的，原文紧接着解释了 $p=1$ 时 $L^1$-界自动成立、因而不提供任何信息。抄写时勿把绝对值与 $p$ 次幂写到分母上。]

\begin{remark}{5.6}
要体会 Kolodziej 结果的强度，应把它们与至少三维的紧黎曼流形上的 Poisson 方程 $\Delta u=f$ 相对照。若 $f\in L^p$（某个 $p>1$），我们只能推出 $u\in W^{2,p}$，而当 $p$ 接近 $1$ 时 $W^{2,p}$ 不能嵌入 $L^\infty$，因此不能期望 $u$ 的任何先验 $L^\infty$-界。归根结底，是整体正性条件 $\phi\in PSH(X,\omega)$ 造成了差别。
\end{remark}

\begin{remark}{5.7}
Kolodziej 的证明依赖于他关于复势理论的“容量衰减”论证。作者从 Freid Tong 处得知，Kolodziej 的结果中有相当一部分已经获得新的证明 [69][70]，其灵感来自 Chen 与 Cheng [15] 关于常数量曲率 Kähler 方程的最新突破。
\end{remark}

% [Tutor Note: 注 5.6 中的对比对象是“Poisson 方程”而不是 Monge--Ampère 方程；原文用这个对比说明：一般的椭圆方程在 $L^p$ 密度假设下得不到 $L^\infty$ 界，而复 Monge--Ampère 方程因势的整体正性（$\phi\in PSH$）而具有额外的刚性。$W^{2,p}\subset L^\infty$ 的失效只在 $p$ 接近 $1$（且维数 $\ge3$）时发生。]

在我们的应用中，我们需要处理一族 Kähler 流形，而估计必须在非常剧烈的复结构退化下保持一致，并允许总体积坍缩到零。一些微妙之处是：
